% !TEX root = ../main.tex
%% Chapitre 19 — Procédés sommatoires discrets
%% Programme officiel MPSI/MP2I (BO spécial n°1 du 11/02/2021), deuxième semestre :
%% a) séries numériques : convergence, divergence grossière, restes, lien
%% suite-série, géométriques, e^z ; b) séries à termes positifs (comparaison,
%% équivalents, série-intégrale, Riemann) ; c) convergence absolue ;
%% d) théorème des séries alternées ; e) familles sommables positives
%% (paquets, Fubini positif) ; f) familles sommables complexes (l¹, Fubini,
%% produit de Cauchy, exponentielle morphisme). Critère de Cauchy hors programme.
%% Résultats vérifiés par verifications/chapitre_19.py.

\begin{ouverture}[sommes]{19}{Procédés sommatoires discrets}{Partie 2 · Deuxième semestre}
  {Séries\enspace·\enspace Comparaisons\enspace·\enspace Riemann\enspace·\enspace
   Séries alternées\enspace·\enspace Familles sommables\enspace·\enspace Fubini}
  \objectif{Étudier la convergence d'une série : divergence grossière, télescopages, géométriques.}
  \objectif{Comparer les séries à termes positifs : inégalités, équivalents, série-intégrale, Riemann.}
  \objectif{Utiliser la convergence absolue et le théorème des séries alternées.}
  \objectif{Sommer « en vrac » : familles sommables, sommation par paquets, Fubini.}
  \objectif{Multiplier des séries : le produit de Cauchy, et l'exponentielle morphisme.}
  \objectif{Savoir quand l'ordre des termes compte — et quand il ne compte plus.}
  \prerequis{Suites : limites, monotonie, encadrements (chapitre 5).}
  \prerequis{Analyse asymptotique : équivalents, développements (chapitre 11).}
  \prerequis{Comparaison série-intégrale (chapitre 15) ; nombres complexes (chapitre 3).}
\end{ouverture}

%% =====================================================================
%%  L'ESSENTIEL
%% =====================================================================
\essentiel

\essSection{1}{Séries numériques : le vocabulaire}

\begin{definition}[Convergence d'une série]
La \textbf{série} $\sum u_n$ est la suite de ses \emph{sommes partielles}
$S_N=\sum_{n=0}^{N}u_n$. Elle \emph{converge} si $(S_N)$ converge ; sa
\textbf{somme} est alors $\sum_{n=0}^{+\infty}u_n$, et son \emph{reste}
$R_N=\sum_{n=N+1}^{+\infty}u_n$ tend vers $0$. La somme est linéaire. Le
terme général d'une série convergente tend vers $0$ : sinon, la série
\emph{diverge grossièrement} — premier réflexe de l'étude.
\end{definition}

\begin{theoreme}[Suite-série ; géométriques ; exponentielle]
La suite $(u_n)$ et la série télescopique $\sum(u_{n+1}-u_n)$ sont de même
nature : les séries contiennent les suites. La série géométrique
$\sum z^{n}$ converge si et seulement si $|z|<1$, de somme
$\dfrac{1}{1-z}$. Et pour tout $z\in\Cx$ :
\[
\mathrm{e}^{z}=\sum_{n=0}^{+\infty}\frac{z^{n}}{n!}
\]
— la série exponentielle, absolument convergente partout.
\end{theoreme}

\essSection{2}{Séries à termes positifs}

\noindent
\begin{minipage}[c]{0.55\linewidth}
\begin{theoreme}[Comparaisons]
Une série à termes positifs converge si et seulement si ses sommes
partielles sont \emph{majorées}. Si $0\leqslant u_n\leqslant v_n$ : la
convergence de $\sum v_n$ entraîne celle de $\sum u_n$. Si
$u_n\sim v_n$ avec $(u_n),(v_n)$ positives : $\sum u_n$ et $\sum v_n$ sont
de \emph{même nature}. Enfin, pour $f$ positive décroissante, les
rectangles encadrent :
\[
\int_{1}^{N+1}f\leqslant\sum_{n=1}^{N}f(n)\leqslant f(1)+\int_{1}^{N}f ,
\]
d'où l'étude des sommes partielles par les intégrales (chapitre 15).
\end{theoreme}
\end{minipage}\hfill
\begin{minipage}[c]{0.41\linewidth}
\centering
\begin{tikzpicture}[scale=0.6,line cap=round]
  \foreach \k in {1,...,6}{%
    \fill[otGreen,opacity=0.28] (\k,0) rectangle ({\k+1},{2.2/\k});
    \draw[otGreenDark,line width=0.5pt] (\k,0) rectangle ({\k+1},{2.2/\k});}
  \draw[otMuted,line width=0.5pt,-{Stealth[length=1.6mm]}] (0.4,0) -- (7.6,0);
  \draw[otMuted,line width=0.5pt,-{Stealth[length=1.6mm]}] (0.7,-0.25) -- (0.7,2.9);
  \draw[otNavy,line width=1.2pt]
    plot[domain=0.85:7.4,samples=70] ({\x},{2.2/\x});
  \node[font=\scriptsize,text=otNavy] at (6.4,0.85) {$y=f(x)$};
\end{tikzpicture}\\[0.5mm]
\legende{Série contre intégrale : les rectangles de hauteur $f(n)$
encadrent l'aire sous la courbe — et réciproquement.}
\end{minipage}

\begin{theoreme}[Séries de Riemann]
$\displaystyle\sum\frac{1}{n^{\alpha}}$ converge si et seulement si
$\alpha>1$ (comparaison série-intégrale avec $t\mapsto t^{-\alpha}$). La
frontière est la série harmonique $\sum\frac1n$, divergente en
$\ln n$ — et $\sum\frac{1}{n^{2}}$ converge, pierre angulaire de bien des
comparaisons (dont la preuve de Stirling, chapitre 15).
\end{theoreme}

\essSection{3}{Convergence absolue}

\begin{theoreme}[ACV $\Rightarrow$ CV]
Une série numérique \textbf{absolument convergente}
($\sum|u_n|$ converge) est convergente. En pratique : si
$u_n=O(v_n)$ avec $(v_n)$ positive sommable, alors $\sum u_n$ converge
absolument. (Le critère de Cauchy est hors programme.) La réciproque est
fausse : $\sum\frac{(-1)^{n-1}}{n}$ converge (vers $\ln2$) sans converger
absolument — elle est \emph{semi-convergente}, et l'ordre de ses termes
devient alors essentiel (exercice 15).
\end{theoreme}

\essSection{4}{Le théorème des séries alternées}

\begin{theoreme}[Séries alternées]
Si $(u_n)$ décroît vers $0$, alors $\sum(-1)^{n}u_n$ converge. De plus, la
somme est comprise entre deux sommes partielles consécutives, chaque reste
est du signe de son premier terme, et
\[
|R_N|\leqslant u_{N+1} :
\]
l'erreur est contrôlée par le premier terme négligé. Exemples :
$\sum\frac{(-1)^{n-1}}{n}=\ln2$ et $\sum\frac{(-1)^{n}}{2n+1}=\frac\pi4$
— déjà rencontrées aux chapitres 11 et 15 par les restes intégraux.
\end{theoreme}

\essSection{5}{Familles sommables de réels positifs}

\begin{definition}[Sommer en vrac]
Dans $\intFF{0}{+\infty}$ (muni de ses conventions de calcul), la somme
d'une famille $(u_i)_{i\in I}$ de réels \emph{positifs} est
\[
\sum_{i\in I}u_i=\sup\Bigl\{\sum_{i\in F}u_i\ \Bigm|\ F\subseteq I
\text{ finie}\Bigr\}\in\intFF{0}{+\infty} :
\]
aucun ordre n'est requis. La famille est \textbf{sommable} si cette somme
est finie. Pour $I=\N$, on retrouve les séries à termes positifs ; la
somme est invariante par permutation.
\end{definition}

\begin{theoreme}[Paquets et Fubini positifs]
Si $I$ est réunion disjointe des $I_j$ :
$\displaystyle\sum_{i\in I}u_i=\sum_{j\in J}\Bigl(\sum_{i\in I_j}u_i\Bigr)$
(démonstration hors programme) — on peut regrouper comme on veut. Cas
d'un produit $I\times I'$ : le théorème de \emph{Fubini positif} autorise
à sommer dans les deux ordres, lignes ou colonnes. La positivité justifie
\emph{tous} les calculs — et leur résultat (fini ou non) décide de la
sommabilité.
\end{theoreme}

\essSection{6}{Familles sommables complexes ; produit de Cauchy}

\begin{theoreme}[Sommabilité complexe]
Une famille $(u_i)_{i\in I}$ de complexes est \textbf{sommable}
(notation $\ell^{1}(I)$) si $\sum_{i\in I}|u_i|<+\infty$ ; sa somme est
alors bien définie, linéaire, invariante par permutation, et toute
sous-famille est sommable. Critère de domination : si
$|u_i|\leqslant v_i$ avec $(v_i)$ sommable positive, alors
$(u_i)$ est sommable. Les théorèmes de sommation par paquets et de
Fubini s'étendent aux familles sommables, ainsi que le produit de deux
familles sommables.
\end{theoreme}

\begin{theoreme}[Produit de Cauchy]
Si $\sum a_n$ et $\sum b_n$ convergent absolument, leur \textbf{produit
de Cauchy} $\sum c_n$, avec $c_n=\sum_{k=0}^{n}a_kb_{n-k}$, converge
absolument et
\[
\sum_{n=0}^{+\infty}c_n
=\Bigl(\sum_{n=0}^{+\infty}a_n\Bigr)\Bigl(\sum_{n=0}^{+\infty}b_n\Bigr) .
\]
Appliqué aux séries exponentielles, le binôme donne
$\mathrm{e}^{a}\mathrm{e}^{b}=\mathrm{e}^{a+b}$ : \emph{l'exponentielle
complexe est un morphisme de $(\Cx,+)$ dans $(\Cx^{*},\times)$} —
l'identité fondatrice du chapitre 3, enfin démontrée.
\end{theoreme}

\begin{attention}
Hors sommabilité, tout s'effondre : une série semi-convergente
\emph{réarrangée} peut changer de somme (exercice 15), et un tableau
double non sommable peut donner deux totaux différents selon l'ordre de
sommation (exercice 24). Sommer en vrac est un droit qui se paie en
valeur absolue.
\end{attention}

%% =====================================================================
%%  EXERCICES
%% =====================================================================
\exercices
\rubrique{Applications}

\begin{exercice}[1]{Convergence ou divergence ?}
Étudier la nature des séries :
\begin{tasks}(2)
  \task $\displaystyle\sum\Bigl(\frac12\Bigr)^{n}$
  \task $\displaystyle\sum\frac{n}{n+1}$
  \task $\displaystyle\sum(-1)^{n}$
  \task $\displaystyle\sum\frac{1}{n(n+1)}$
\end{tasks}
Donner les sommes quand elles existent.
\end{exercice}

\begin{exercice}[1]{Séries géométriques et exponentielle}
Calculer :
\begin{tasks}(3)
  \task $\displaystyle\sum_{n=0}^{+\infty}\frac{1}{3^{n}}$
  \task $\displaystyle\sum_{n=0}^{+\infty}\frac{(-1)^{n}}{2^{n}}$
  \task $\displaystyle\sum_{n=0}^{+\infty}\frac{1}{n!}$
\end{tasks}
\end{exercice}

\begin{exercice}[1]{Télescopages}
\begin{questions}
  \item Calculer $\displaystyle\sum_{n=1}^{+\infty}\frac{1}{n(n+1)}$ puis
    $\displaystyle\sum_{n=1}^{+\infty}\frac{1}{n(n+1)(n+2)}$.
  \item Montrer que $\displaystyle\sum\ln\Bigl(1+\frac1n\Bigr)$ diverge,
    en calculant ses sommes partielles.
\end{questions}
\end{exercice}

\begin{exercice}[1]{Comparer des séries positives}
\begin{questions}
  \item Montrer que $\displaystyle\sum\frac{1}{n^{2}}$ converge en
    comparant à $\dfrac{1}{n(n-1)}$.
  \item Nature de $\displaystyle\sum\frac{1}{n^{2}+n+1}$ et de
    $\displaystyle\sum\frac{1}{\sqrt n}$ ?
\end{questions}
\end{exercice}

\begin{exercice}[1]{Le critère des équivalents}
Déterminer la nature de :
\begin{tasks}(2)
  \task $\displaystyle\sum\sin\frac{1}{n^{2}}$
  \task $\displaystyle\sum\Bigl(1-\cos\frac1n\Bigr)$
  \task $\displaystyle\sum\ln\Bigl(1+\frac1n\Bigr)$
  \task $\displaystyle\sum\frac{\sqrt{n+1}-\sqrt n}{n}$
\end{tasks}
\end{exercice}

\begin{exercice}[1]{Les séries de Riemann}
\begin{questions}
  \item Démontrer le critère : $\displaystyle\sum\frac{1}{n^{\alpha}}$
    converge si et seulement si $\alpha>1$ (comparaison
    série-intégrale).
  \item Appliquer à $\alpha=2$, $\alpha=\dfrac32$, $\alpha=\dfrac12$.
\end{questions}
\end{exercice}

\begin{exercice}[1]{Sommes partielles par série-intégrale}
\begin{questions}
  \item Montrer que
    $\displaystyle\sum_{k=1}^{n}\frac{1}{\sqrt k}\sim2\sqrt n$.
  \item Retrouver $H_n\sim\ln n$, et le lien avec la constante d'Euler
    (chapitre 15).
\end{questions}
\end{exercice}

\begin{exercice}[1]{Séries alternées : premiers pas}
\begin{questions}
  \item Montrer que $\displaystyle\sum\frac{(-1)^{n-1}}{n}$ et
    $\displaystyle\sum\frac{(-1)^{n}}{\sqrt n}$ convergent.
  \item Combien de termes de la première garantissent une erreur
    inférieure à $10^{-2}$ ? Quelle est sa somme (chapitres 11 et 15) ?
\end{questions}
\end{exercice}

\rubrique{Entraînement}

\begin{exercice}[2]{Du lien suite-série}
On pose $u_n=\displaystyle\prod_{k=1}^{n}\Bigl(1+\frac{1}{k^{2}}\Bigr)$.
\begin{questions}
  \item Montrer que $(u_n)$ converge, en étudiant la série
    $\sum\ln\Bigl(1+\dfrac{1}{k^{2}}\Bigr)$.
  \item En déduire que les produits infinis se ramènent aux séries.
\end{questions}
\end{exercice}

\begin{exercice}[2]{Restes de séries convergentes}
\begin{questions}
  \item Encadrer $R_N=\displaystyle\sum_{n=N+1}^{+\infty}\frac{1}{n^{2}}$
    entre $\dfrac{1}{N+1}$ et $\dfrac1N$, et en déduire un équivalent.
  \item Calculer exactement le reste de la série
    $\sum\bigl(\tfrac12\bigr)^{n}$ : quelle vitesse, comparée à a) ?
\end{questions}
\end{exercice}

\begin{exercice}[2]{Convergence absolue}
\begin{questions}
  \item Montrer que $\displaystyle\sum\frac{\sin n}{n^{2}}$ converge.
  \item La série harmonique alternée converge-t-elle absolument ?
    Qu'appelle-t-on série semi-convergente ?
\end{questions}
\end{exercice}

\begin{exercice}[2]{Le produit de Cauchy et l'exponentielle}
\begin{questions}
  \item Former le produit de Cauchy des séries $\sum\dfrac{a^{n}}{n!}$ et
    $\sum\dfrac{b^{n}}{n!}$, et reconnaître le binôme : en déduire
    $\mathrm{e}^{a}\mathrm{e}^{b}=\mathrm{e}^{a+b}$.
  \item Former le produit de Cauchy de $\sum x^{n}$ par elle-même, et en
    déduire $\displaystyle\sum_{n=0}^{+\infty}(n+1)x^{n}
    =\frac{1}{(1-x)^{2}}$ pour $|x|<1$.
\end{questions}
\end{exercice}

\begin{exercice}[2]{L'exponentielle complexe, enfin fondée}
\begin{questions}
  \item Justifier la convergence absolue de
    $\sum\dfrac{z^{n}}{n!}$ pour tout $z\in\Cx$, et l'inégalité
    $\bigl|\mathrm{e}^{z}\bigr|\leqslant\mathrm{e}^{|z|}$.
  \item En séparant parties réelle et imaginaire de
    $\mathrm{e}^{\mathrm{i}\theta}$, retrouver les séries de
    $\cos\theta$ et $\sin\theta$, et le cercle trigonométrique du
    chapitre 3.
\end{questions}
\end{exercice}

\begin{exercice}[2]{Sommes du type $\sum n x^{n}$}
\begin{questions}
  \item Établir $\displaystyle\sum_{n=1}^{+\infty}nx^{n}
    =\frac{x}{(1-x)^{2}}$ pour $|x|<1$ (exercice 12, ou télescopage).
  \item En déduire $\displaystyle\sum_{n=1}^{+\infty}\frac{n}{2^{n}}=2$,
    puis calculer $\displaystyle\sum_{n=1}^{+\infty}\frac{n^{2}}{2^{n}}$.
\end{questions}
\end{exercice}

\begin{exercice}[2]{Le réarrangement qui change la somme}
On réordonne la série harmonique alternée en prenant \emph{deux} termes
positifs puis \emph{un} négatif :
$1+\dfrac13-\dfrac12+\dfrac15+\dfrac17-\dfrac14+\cdots$
\begin{questions}
  \item En regroupant par blocs de trois, montrer que la nouvelle série
    converge vers $\dfrac32\ln2$ — et non $\ln2$.
  \item Pourquoi ce phénomène est-il impossible pour une série
    absolument convergente ?
\end{questions}
\end{exercice}

\begin{exercice}[2]{Fubini positif : premiers calculs}
\begin{questions}
  \item Calculer $\displaystyle\sum_{(i,j)\in\N^{2}}\frac{1}{2^{i+j}}$ de
    deux façons : comme produit de deux géométriques, puis par paquets
    diagonaux $i+j=n$.
  \item Que donne la cohérence des deux calculs ?
\end{questions}
\end{exercice}

\begin{exercice}[2]{Sommation par paquets}
Calculer $\displaystyle\sum_{n=2}^{+\infty}\frac{1}{n^{2}-1}$ en
décomposant en éléments simples (chapitre 10) et en regroupant les termes
télescopiques. Pourquoi le regroupement est-il licite ?
\end{exercice}

\begin{exercice}[2]{Fubini en action}
Pour $0<x<1$, calculer de deux façons la somme de la famille
$\bigl(x^{j}\bigr)_{(i,j),\ j\geqslant i\geqslant0}$ :
\begin{questions}
  \item à $i$ fixé, puis en sommant sur $i$ ;
  \item à $j$ fixé, puis en sommant sur $j$.
\end{questions}
En déduire une nouvelle preuve de
$\sum_{j\geqslant0}(j+1)x^{j}=\dfrac{1}{(1-x)^{2}}$.
\end{exercice}

\rubrique{Approfondissement}

\begin{exercice}[3]{Encadrer les valeurs de $\zeta$}
On note $\zeta(\alpha)=\displaystyle\sum_{n=1}^{+\infty}\frac{1}{n^{\alpha}}$
pour $\alpha>1$.
\begin{questions}
  \item Par série-intégrale, montrer que
    $\dfrac32\leqslant\zeta(2)\leqslant2$ et
    $\dfrac98\leqslant\zeta(3)\leqslant\dfrac32$.
  \item Euler a montré $\zeta(2)=\dfrac{\pi^{2}}{6}\approx1{,}645$ (hors
    programme) : vérifier la cohérence avec a).
\end{questions}
\end{exercice}

\begin{exercice}[3]{Développement asymptotique d'un reste}
Soit $R_N=\displaystyle\sum_{n=N+1}^{+\infty}\frac{1}{n^{2}}$.
\begin{questions}
  \item Encadrer $\dfrac{1}{n^{2}}$ entre $\dfrac{1}{n(n+1)}$ et
    $\dfrac{1}{(n-1)n}$ et en déduire
    $\dfrac{1}{N+1}\leqslant R_N\leqslant\dfrac1N$, puis
    $R_N\sim\dfrac1N$.
  \item Avec le télescopage médian
    $\dfrac{1}{n-\frac12}-\dfrac{1}{n+\frac12}=\dfrac{1}{n^{2}-\frac14}
    >\dfrac{1}{n^{2}}$, affiner : $R_N<\dfrac{1}{N+\frac12}$, et situer
    $R_N=\dfrac1N+O\Bigl(\dfrac{1}{N^{2}}\Bigr)$ (chapitre 11).
\end{questions}
\end{exercice}

\begin{exercice}[3]{La constante d'Euler par les séries}
\begin{questions}
  \item Montrer que la série
    $\displaystyle\sum_{k\geqslant1}\Bigl(\frac1k-\ln\Bigl(1+\frac1k\Bigr)\Bigr)$
    converge. \emph{Équivalent du terme général.}
  \item Montrer que sa somme est la constante d'Euler $\gamma$
    (chapitre 15) : $\gamma$ est la somme d'une série à termes positifs
    équivalents à $\frac{1}{2k^{2}}$.
\end{questions}
\end{exercice}

\begin{exercice}[3]{Un produit eulérien miniature}
Soit $A$ l'ensemble des entiers de la forme $2^{a}3^{b}$
($a,b\in\N$).
\begin{questions}
  \item Montrer que la famille $\bigl(\frac1n\bigr)_{n\in A}$ est
    sommable, de somme
    $\dfrac{1}{1-\frac12}\cdot\dfrac{1}{1-\frac13}=3$. \emph{Fubini
    positif sur les exposants.}
  \item Pourquoi la même idée, étendue à tous les nombres premiers,
    relie-t-elle la série harmonique à l'infinité des premiers
    (chapitre 7) ?
\end{questions}
\end{exercice}

\begin{exercice}[3]{Le critère de condensation de Cauchy}
Soit $(u_n)$ positive décroissante.
\begin{questions}
  \item Par sommation par paquets dyadiques
    $\ens{2^{k},\dots,2^{k+1}-1}$, montrer que $\sum u_n$ converge si et
    seulement si $\sum2^{k}u_{2^{k}}$ converge.
  \item Retrouver le critère de Riemann, puis montrer que
    $\displaystyle\sum\frac{1}{n\ln n}$ diverge — plus lentement encore
    que l'harmonique.
\end{questions}
\end{exercice}

\begin{exercice}[3]{Quand Fubini s'effondre}
On considère le tableau $(a_{ij})_{i,j\in\N}$ défini par
$a_{ii}=1$, $a_{i+1,i}=-1$, et $0$ ailleurs.
\begin{questions}
  \item Calculer la somme de chaque ligne, puis la somme de ces sommes ;
    faire de même avec les colonnes. Que constate-t-on ?
  \item Pourquoi le théorème de Fubini ne s'applique-t-il pas ? Vérifier
    que la famille $\bigl(|a_{ij}|\bigr)$ n'est pas sommable.
  \item Conclure : la sommabilité n'est pas une précaution de juriste,
    mais la condition exacte qui rend l'ordre de sommation indifférent —
    le mot de la fin du chapitre.
\end{questions}
\end{exercice}

% !TEX root = ../../main.tex
%% Défi du chapitre 19 — Le problème de Bâle
%% Résultats vérifiés par verifications/defis_16_20.py.

\rubrique{Défi}

\begin{defi}{Le problème de Bâle : $\zeta(2)=\pi^{2}/6$}{2 h}
Posé en 1644 par Pietro Mengoli, résolu par Leonhard Euler en 1735, le
« problème de Bâle » demande la valeur exacte de
$\zeta(2)=\displaystyle\sum_{k=1}^{+\infty}\frac{1}{k^{2}}$. On en donne une
preuve entièrement élémentaire, par la trigonométrie et les relations de
Viète. Soit $m\in\N^{*}$.

\partieExo{A}{Une identité polynomiale}
\begin{questions}
  \item À l'aide de la formule de Moivre appliquée à
    $(\cos\theta+\mathrm{i}\sin\theta)^{2m+1}$, montrer que
    \[
    \sin\bigl((2m+1)\theta\bigr)=\sum_{k=0}^{m}(-1)^{k}\binom{2m+1}{2k+1}
    \cos^{2m-2k}\theta\,\sin^{2k+1}\theta .
    \]
  \item En déduire, pour $\theta\in\intOO{0}{\frac\pi2}$ :
    $\dfrac{\sin\bigl((2m+1)\theta\bigr)}{\sin^{2m+1}\theta}=P_m\bigl(\cot^{2}\theta\bigr)$,
    où $P_m=\sum_{k=0}^{m}(-1)^{k}\binom{2m+1}{2k+1}X^{m-k}$ et
    $\cot=\frac{\cos}{\sin}$.
  \item Montrer que les $m$ réels $x_r=\cot^{2}\frac{r\pi}{2m+1}$
    ($1\leqslant r\leqslant m$) sont des racines distinctes de $P_m$ ; en
    déduire que ce sont toutes ses racines, puis, par les relations de
    Viète :
    \[
    \sum_{r=1}^{m}\cot^{2}\frac{r\pi}{2m+1}=\frac{m(2m-1)}{3} .
    \]
\end{questions}

\partieExo{B}{L'encadrement}
\begin{questions}[resume]
  \item Montrer que $\sin\theta<\theta<\tan\theta$ pour
    $\theta\in\intOO{0}{\frac\pi2}$, et en déduire
    $\cot^{2}\theta<\dfrac{1}{\theta^{2}}<1+\cot^{2}\theta$.
  \item En sommant ces inégalités pour $\theta=\frac{r\pi}{2m+1}$,
    $1\leqslant r\leqslant m$, montrer que
    \[
    \frac{m(2m-1)}{3}<\frac{(2m+1)^{2}}{\pi^{2}}\sum_{r=1}^{m}\frac{1}{r^{2}}
    <\frac{2m(m+1)}{3} .
    \]
  \item En déduire que $\zeta(2)=\dfrac{\pi^{2}}{6}$.
\end{questions}

\medskip
\begin{center}
\begin{tikzpicture}[x=3cm,y=0.42cm,line cap=round]
  \foreach \v in {2,4,6,8}{\draw[otRule,line width=0.3pt] (0,\v) -- (1.62,\v);
    \node[font=\tiny,text=otMuted,anchor=east] at (0,\v) {$\v$};}
  \draw[otMuted,line width=0.5pt,-{Stealth[length=1.6mm]}] (0,0) -- (1.75,0);
  \draw[otMuted,line width=0.5pt,-{Stealth[length=1.6mm]}] (0,0) -- (0,9.2);
  \draw[otMuted,line width=0.4pt,dash pattern=on 2pt off 1.5pt] (1.5708,0) -- (1.5708,8.6);
  \draw[otNavy!45,line width=1pt] plot[domain=0.343:1.5708,samples=120] ({\x},{1/(sin(deg(\x))^2)});
  \draw[otGreenDark,line width=1.2pt] plot[domain=0.343:1.5708,samples=120] ({\x},{1/(\x*\x)});
  \draw[otNavy,line width=1pt] plot[domain=0.36:1.5708,samples=120] ({\x},{(cos(deg(\x))/sin(deg(\x)))^2});
  \node[font=\tiny,text=otMuted,below] at (1.5708,0) {$\frac\pi2$};
  \node[font=\scriptsize,text=otNavy!60,anchor=west] at (0.62,3.6) {$1+\cot^{2}\theta$};
  \node[font=\scriptsize,text=otGreenDark,anchor=west] at (0.86,1.9) {$1/\theta^{2}$};
  \node[font=\scriptsize,text=otNavy,anchor=north east] at (0.78,0.7) {$\cot^{2}\theta$};
\end{tikzpicture}\\[0.5mm]
\legende{Sur $\intOO{0}{\frac\pi2}$, la courbe de $1/\theta^{2}$ (en vert) est
coincée entre celles de $\cot^{2}\theta$ et de $1+\cot^{2}\theta=1/\sin^{2}\theta$.}
\end{center}

\partieExo{C}{Conséquences}
\begin{questions}[resume]
  \item En séparant termes pairs et impairs, en déduire
    $\displaystyle\sum_{k\geqslant0}\frac{1}{(2k+1)^{2}}=\frac{\pi^{2}}{8}$ et
    $\displaystyle\sum_{k\geqslant1}\frac{(-1)^{k+1}}{k^{2}}=\frac{\pi^{2}}{12}$.
  \item On admet l'encadrement du reste
    $\frac{1}{N+1}<\sum_{k>N}\frac{1}{k^{2}}<\frac{1}{N+\frac12}$
    (chapitre 19, exercice 20). Sachant que
    $\sum_{k=1}^{10}\frac1{k^{2}}\approx1{,}549\,768$, en déduire un
    encadrement de $\pi$ d'amplitude inférieure à $5\cdot10^{-3}$.
\end{questions}
\end{defi}

